\section{A Pell descent and a degree-82 equation in twelve positive unknowns}
\label{sec:pell-saturation}

The principal new equation in this component has twelve strictly positive
integer unknowns, exact total degree $82$, and $102$ live binary arithmetic
operations.  The inspected package61 equation has the same domain and
variable count, degree $118$, and $104$ operations.  Thus this construction
reduces its degree by $30.5085\%$ and saves two operations.  The fixed
coefficient tuple is unchanged.  The proof uses an elementary descent to
show that a rational auxiliary Pell ordinate is automatically integral;
this permits two source witnesses to be merged into one.

\subsection{Rational ordinates of a particular Pell conic are integral}

\paragraph{Lemma.}
Let $Q>1$ be an integer.  If $V\in\mathbb Z$, $y\in\mathbb Q$, and
\begin{equation}
 QV^2-(Q-1)y^2=1,                                      \label{eq:sat-pell}
\end{equation}
then $y\in\mathbb Z$.

\paragraph{Proof.}
Neither coordinate vanishes, so sign changes let us take $V,y>0$.
Write $y=a/m$ in lowest terms, $m>0$.  Equation~\eqref{eq:sat-pell}
implies $m^2\mid Q-1$.  If $V=1$, then $y=1$ and $m=1$.
For $V>1$, set
\begin{align}
 V'&=(2Q-1)V-2(Q-1)y,\nonumber\\
 y'&=(2Q-1)y-2QV.                                      \label{eq:sat-step}
\end{align}
The first coordinate is an integer.  Since
$(2Q-1)^2-4Q(Q-1)=1$, direct expansion gives
$Q(V')^2-(Q-1)(y')^2=1$.  Furthermore
\[
 (Q-1)(y^2-V^2)=V^2-1>0,
 \qquad
 (QV)^2-((Q-1)y)^2=QV^2+Q-1>0.
\]
Thus $y>V$ and $QV>(Q-1)y$, so $-V<V'<V$.
The norm excludes $V'=0$.  Finally $y'$ has exactly the same reduced
denominator $m$, because $2Q-1\equiv1\pmod m$ and the other summand is
integral.  Strong induction on the positive integer $V$, replacing the
new coordinates by their absolute values, reaches $V=1$ without changing
$m$.  Therefore $m=1$.\hfill$\square$

The adjacent coefficients $Q,Q-1$ are essential to this proof.
For example $2\cdot1^2-4\cdot(1/2)^2=1$ has a nonintegral rational
ordinate.  No coprimality condition involving another equation is needed
in the displayed lemma.

\subsection{The exact norm merger}

For positive integers $d,\rho,Q,V,n$ with $Q>1$, define
\begin{align}
 h(z)&=(1-Q)z^2+2(1-Q)Vz+V^2-1,\nonumber\\
 a_N&=h(n),& b_N&=2(1-Q)(n+V),\nonumber\\
 A_N&=d a_N+(1-Q)\rho,&
 N&=A_N^2-d b_N^2\rho.                                 \label{eq:sat-norm}
\end{align}
Here $b_N\ne0$.  If $N=0$, the proof-only rational numbers
\[
 i=\frac{A_N}{d b_N},\qquad u=n-i
\]
satisfy $di^2=\rho$ and $h(u)=0$.  The second equality is
$QV^2-(Q-1)(V+u)^2=1$.  The lemma makes $V+u$ integral; consequently
$u$ and $i$ are integral as well.  Conversely $di^2=\rho$ and
$h(u)=0$, with $n=i+u$, imply $N=0$.  The exact identity is
\begin{equation}
 N\big|_{\rho=di^2}=d^2h(n-i)h(n+i).                    \label{eq:sat-factor}
\end{equation}
The norm identity by itself gives rational coordinates.  The descent is
the step that recovers integer coordinates.
If $V>1$, the auxiliary norm gives $|V+u|>|V|$.  Thus a positive lift is
obtained by replacing $i$ with $|i|$ and $u$ with $|V+u|-V$.

\subsection{The complete polynomial}

Supply exactly
\begin{equation}
 j,F,\alpha,t,T,X,Y,r,\eta,\zeta,n,\gamma\in\mathbb Z_{>0}.
                                                               \label{eq:sat-vars}
\end{equation}
All subsequent symbols are polynomial abbreviations, not extra unknowns.
The four fixed numerals are the complete compact44 tuple
$\beta,K_0,M_C,M_F$, where $\beta=2^{590069}-1$ and $M_F$ is already
the shifted high mask.  Define the retained carrier registers by
\begin{align}
 J&=6j+1,& q&=\beta J+1,& C&=q-F-\alpha,\nonumber\\
 E&=XY,& k&=\eta+\zeta,& c&=kY+\eta,\nonumber\\
 a&=E+Y,& H&=4a+3,& \Delta&=a^2+H=(a+2)^2-1,\nonumber\\
 \ell&=EkY,& b_0&=X+ac,\nonumber\\
 R&=((q-1)(q-F)+\alpha)(q^2-1)+(M_C+qM_F)J,\nonumber\\
 v&=k-R-1,& U&=q^5,\nonumber\\
 D_*&=Ub_0+H(U\gamma+q^4X-v),\nonumber\\
 W&=(q-1)k+1,& M&=q^3(q-1),\nonumber\\
 B_t&=(qK_0+X)C-qF,\nonumber\\
 Z&=M(\ell+k+r)+(q-1)Y-q^3WB_t,\nonumber\\
 A_0&=Z^2-M^2[\ell(\ell+k)+1],\nonumber\\
 A_1&=D_*^2-U^2(1+\Delta c^2).                         \label{eq:sat-carriers}
\end{align}
For the new norm put
\begin{equation}
 f=t+1,\quad \rho=f^2-1,\quad d=\Delta c^4,
 \quad Q=\Delta\rho,\quad V=f^2(cT-R)-c,                 \label{eq:sat-inner}
\end{equation}
and use~\eqref{eq:sat-norm}.  The final equation is
\begin{equation}
 \boxed{P_{12}=(1+A_0)(1+A_1)(1+N)-1=0.}                \label{eq:sat-P12}
\end{equation}
The literal evaluator shares $s=V+n$ and computes
$h(n)=s^2-Q(s^2-V^2)-1$.  Every operation is $+$, $-$, or
multiplication; every exponent is fixed.

\subsection{Soundness and the order of the inherited bounds}

Since $J$ and $\beta$ are odd, $q$ is even.  Thus $A_0$ and $A_1$
are squares modulo four.  The same is true of $N$ because $b_N$ is even.
If~\eqref{eq:sat-P12} vanishes, each integer factor is $+1$ or $-1$.
A factor $-1$ would require a residual $-2\equiv2\pmod4$, impossible
for a square modulo four.  Therefore $A_0=A_1=N=0$.

Package61's retained carrier argument, before it invokes either strong
or auxiliary Pell equation, proves the following.  From $A_0=0$ it
recovers $q^3\mid Y$, a positive transport quotient, and $C>0$.
The authentic coefficient bounds then imply
\[
 q\ge16,\qquad 3q<R<q^4-2.
\]
The second carrier gives a positive main root and the first carrier a
possibly signed first root.  Positive Pell classification gives indices
$p,m_1\ge1$ with
\[
 c=\psi_{a+2}(p),\qquad
 k=2\psi_{2XY^2+1}(m_1).
\]
The case $m_1=1$ would give $k=2$ and $c=2Y+1$, contradicting
$c\ge2(a+2)\ge4Y+4$.  Hence $m_1\ge2$, and
\begin{equation}
 c>kY\ge4(2XY^2+1)Y>8XY^3\ge8q^9>2R.                 \label{eq:sat-bound}
\end{equation}
This uses only the two carriers; it does not presume the rank recovery
that follows from the strong and auxiliary norms.  Because $f=t+1\ge2$
and $T\ge1$, it now follows that
\[
 V=f^2(cT-R)-c>f^2c/2-c\ge c>1.
\]
The norm merger is therefore nondegenerate and gives positive restored
coordinates $i_+,u_+$ satisfying
\[
 f^2-\Delta(i_+c^2)^2=1,
 \qquad QV^2-(Q-1)(V+u_+)^2=1.
\]
Together with the carriers these are exactly package61's thirteen-positive
equations, with its old $j$ coordinate equal to $2j$.  Its remaining rank
recovery and compiler/refutation calibration apply unchanged.

\subsection{Completeness}

The parent's canonical compact witnesses have $B=2^{590069}$,
$q=B^N$, and $J=1+B+\cdots+B^{N-1}$, where $N>1$ is odd.
As $B$ is even, $J$ is odd; as $B=-1\pmod3$, $J=1\pmod3$.
Thus $j=(J-1)/6$ is a positive integer.  Its normalized strong pair has
$f>1$.  Take $t=f-1>0$ and $n=i+u>0$, retaining the other coordinates.
Then $di^2=\rho$ and $h(u)=0$, so~\eqref{eq:sat-factor} gives $N=0$;
both carriers already vanish.  No new congruence restriction or index
doubling is necessary.

Composing the two existence implications with the inherited source
calibration proves, by ordinary mathematics formalizable in ZFC,
\[
 \neg\exists\boldsymbol{x}\in\mathbb Z_{>0}^{12}\,
 P_{12}(\boldsymbol{x})=0
 \quad\Longleftrightarrow\quad \operatorname{Con}(\mathrm{ZFC}).
\]
The new descent and elimination proofs are self-contained.  The native
compiler and its syntactic ZFC endpoint remain the explicitly inherited
semantic dependencies.  This component supplies neither a new complete
proof-assistant certificate nor an exported object-language ZFC derivation.

\subsection{Exact degree and verification}

Set
\[
 q_0=6\beta j,\quad k_0=\eta+\zeta,\quad c_0=k_0Y,
 \quad\Delta_0=X^2Y^2,\quad R_0=q_0^3(q_0-F),
\]
and $B_{t,0}=(K_0q_0+X)(q_0-F-\alpha)-q_0F$.
The exact carrier cancellations give leading forms
\[
 (A_0)_{15}=-2q_0^8k_0^2XY^2B_{t,0},\qquad
 (A_1)_{17}=8q_0^{10}k_0X^2Y^3\gamma.
\]
The degrees of $\rho,d,Q,V,a_N,b_N,A_N$ are respectively
$2,12,6,6,13,12,25$.  Since $db_N^2\rho$ has degree $38$,
\[
 N_{50}=4\Delta_0^4c_0^8t^8R_0^2n^2.
\]
Only $A_0A_1N$ reaches degree $82$ in the final product, with
\begin{equation}
 (P_{12})_{82}=-64q_0^{24}(q_0-F)^2 k_0^{11}
 X^{11}Y^{21}t^8n^2\gamma B_{t,0}.                     \label{eq:sat-leader}
\end{equation}
In particular the coefficient of
$j^{24}F^2\eta^{11}X^{12}Y^{21}t^8n^2\gamma\alpha$
is $64(6\beta)^{24}>0$, independently of the other three ports.
Thus the degree is exactly $82$ on the supplied fixed coefficient slice.

The standard-library checker independently expands both carriers, verifies
the small norm and finalizer cuts exactly, proves the displayed leading
coefficient, checks all $102$ live gates, authenticates the four coefficient
payloads, and rebuilds both the circuit and coefficient tuple in isolated
directories.  The independent native audit also passes.  Normal and
optimized Python runs agree.  These finite checks support transcription
and finite data; they do not replace the unbounded arguments above.
